\FloatBarrier
\section{本章小结}
\label{sec:dp-computation-chapter-summary}

本章在已核验的结构上完成数值计算。所有任务共用同一套组织：每个局部项恰好归入一个容纳其作用域的袋，按后序由叶到根处理，每个节点合并本地的局部项与子节点交来的表，对内部变量汇总，把以接口为索引的结果交给父节点。按消元次序计算与沿消元森林计算得到同样的中间函数；合并袋、调整归属或改变同类汇总的次序会改变中间表，不改变最终结果。

最小成本问题同时求值和恢复最优配置：合并用加法，汇总取最小，取最小时记下达到值，根选定最优状态后沿前序查回溯表，等价于按消元的逆序回溯；最小值为$+\infty$时报告不可行，并列最优按固定规则选取。带保留变量时，查询结构使根袋覆盖保留集，保留变量从不在非根节点被汇总，根处只对其余变量取最小，输出各固定配置下的最小成本，并可恢复相应的条件最优解。计数和求和把加法换成乘法、取最小换成求和，消息表示下方的可行取值数或权重总和，非连通分量的结果相乘，孤立变量贡献其状态数；它们不需要回溯表。

正确性都来自同一个不变量：消息等于固定接口后下方全部局部项的最优值或累计值。它依赖三项条件：运行交集使接口隔开两侧，顶点覆盖使根处涉及全部变量，唯一归属使每项只计一次；再加上组合与汇总运算的结合律、交换律和分配律。这些规则只保证同一种汇总可以换序，求和与取最小混用时不能交换，四项小表就给出$1$与$0$两个不同结果。

实际代价由袋、接口、各变量的状态数、输出要求与存储安排共同决定。宽度给出袋表$q^{w+1}$的上界，比较结构时应使用实际的状态数之积；求值时间约为各袋配置数乘以逐项工作量之和，结构构造与核验另计；存储按同时保留的消息、回溯表和输出计算峰值，可以用释放和重算换取存储。四环算例在全部任务上与枚举逐项一致。

至此，本部分从局部函数出发，经链和树上的接口递推、一般图上的结构构建，完成了有限状态配置问题的精确计算。
